% Generated draft: requires mathematical and source review.
\documentclass[11pt]{article}
\usepackage[margin=25mm]{geometry}
\usepackage{fontspec}
\setmainfont{Latin Modern Roman}
\usepackage{amsmath,amssymb,mathtools,unicode-math}
\setmathfont{Latin Modern Math}
\usepackage{enumitem,needspace,xurl,hyperref}
\setlength{\parindent}{0pt}
\setlength{\parskip}{5pt}
\setlength{\emergencystretch}{3em}
\newcommand{\N}{\mathbb{N}}
\newcommand{\Z}{\mathbb{Z}}
\newcommand{\Q}{\mathbb{Q}}
\newcommand{\R}{\mathbb{R}}
\newcommand{\C}{\mathbb{C}}
\newcommand{\F}{\mathbb{F}}
\newcommand{\A}{\mathbb{A}}
\newcommand{\PP}{\mathbb{P}}
\newcommand{\E}{\mathbb{E}}
\newcommand{\eps}{\varepsilon}
\newcommand{\dd}{\,\mathrm{d}}
\begin{document}
\section*{UnsolvedMath ID: 30002617}
\textbf{Source problem code:} catalogue entry
\subsection{Definitions and mathematical statement}
Write
\[
\mathcal P_{\mu}=\operatorname{Aut}(F\otimes\overline{\Q}/F)[\mu]
\]
for the group expression appearing in the supplied statement. This is retained as source notation: the record does not define the field \(F\), the tensor-product base, the parameter \(\mu\), or the operation denoted by \([\mu]\). Consequently, this expression does not yet specify a group with a topology and an action suitable for continuous cohomology. The notation \(G(\Q)\) denotes the group of rational points of the intended algebraic group \(G\), whose hypotheses and Shimura datum are also absent.

For a specified variety \(X\) over a field \(k\), an algebraic closure \(\bar k\), and a specified coefficient object \(A\), the notation
\[
R\Gamma_{\mathrm{\acute et}}(X_{\bar k},A)
\]
denotes its geometric \(\mathrm{\acute etale}\) cohomology complex. This explains the type of complex mentioned in the question; it does not select the missing base field or coefficients. The source does not specify whether its intended object is the cohomology of an individual Shimura variety, a tower over varying levels, or another cohomology complex, nor whether ordinary or compactly supported cohomology is intended. The prime \(\ell\), coefficient ring, and relevant derived category are not specified.

\textbf{Source question, with unresolved parameters.} For the Shimura variety and cohomology complex intended in the source, does that complex admit a canonical plectic lift carrying the indicated \(\mathcal P_{\mu}\)-action, compatible with products and with the action of \(G(\Q)\), such that the corresponding group cohomology defines plectic \(\ell\)-adic \(\mathrm{\acute etale}\) cohomology?

Here a precise meaning of ``lift'' requires a target category and a specified operation recovering the original cohomology complex and its existing actions. Neither is supplied. The meaning of ``canonical,'' the product compatibility, the \(G(\Q)\)-action, and the group and topology used for group cohomology are likewise unspecified. No universally quantified conjecture over Shimura data or fields can therefore be extracted without additional hypotheses. The question above preserves the supplied target but remains an incomplete formulation requiring mathematical review.
\subsection{Short English statement}
The question asks whether the intended Shimura-variety cohomology complex has a canonical plectic enhancement, compatible with products and \(G(\Q)\), from which plectic \(\ell\)-adic cohomology can be obtained by group cohomology. The supplied record leaves essential fields, groups, coefficients, categories, and compatibility conditions undefined, so a fully precise conjecture cannot yet be stated.
\subsection{Sources}
\begin{itemize}
\item Local UnsolvedMath registry: supplied record with permanent ID \(30002617\), title ``Canonical Plectic Lifts of Shimura-Variety Cohomology,'' exact source problem code \texttt{catalogue entry}, in the Oberwolfach Reports Open Problems collection. No individual proposer is named; the supplied proposed year is 2014.
\item Supplied report citation: \emph{Algebraische Zahlentheorie (2015)}, Oberwolfach Reports, pp. 1747--1800, DOI \texttt{10.4171/owr/2014/32}. \url{https://doi.org/10.4171/owr/2014/32}. The background describes these reports as documenting mathematical workshops and their open-problem discussions. The report citation is supplied provenance, not an independently checked identification of the precise conjecture.
\item Supplied literature reference: Loeffler--Zerbes, ``plectic structures.'' \url{https://arxiv.org/abs/2211.12078}.
\item Supplied literature reference: Li--Huerta, ``function field plectic conjecture.'' \url{https://arxiv.org/abs/2106.05382}.
\item Supplied literature reference: Nekovar--Scholl, ``plectic conjectures.'' \url{https://webusers.imj-prg.fr/~jan.nekovar/pu/introplec.pdf}.
\item Literature assessment dated \texttt{2026-08-21}: the supplied curation labels the problem ``partially solved'' and states that the canonical plectic lift is proved in important special cases, including degree two under mild assumptions and global function fields, while the general number field conjecture remains unresolved. This is supplied dated curation derived from the OWR record, not an independently verified literature assessment; the background explicitly requests independent verification of time-sensitive status and cited claims.
\item The registry labels the statement \texttt{corrected\_verified} and reports that mathematical content was preserved while notation, typography, or omitted definitions were repaired for publication. These are supplied editorial claims, not mathematical verification of this draft.
\end{itemize}
\end{document}
